\centerline{\bf \Huge Тренировочная олимпиада}
\centerline{\bf \Huge 7 класс}

\problem{1} 
Клетки доски $16\times 16$ покрасили в один из двух цветов: красный или синий. Оказалось, что в любом трёхклеточном уголке красных клеток больше, чем синих. Найдите наибольшее возможное число синих клеток.

\problem{2} 
 Докажите, что при $x, y \geq \dfrac{1}{2}$ выполняется неравенство:
 
\[ x^2y^2 + 2(x+y) \geq 4xy + 1 \]

\problem{3} 
Мишень имеет форму правильного $6-$угольника со стороной 1 метр. Техасский бандит Дикарь Герман произвёл $7$ выстрелов по мишени со своего револьвера. Докажите, что  найдутся две пробоины расстояние между которыми меньше метра.

\problem{4}
Мияска рисовала на холсте свой очередной шедевр и случайно разлила краски на лежащую рядом шахматную доску. Некоторые \textit{покрашенные} клетки сразу высохли, а некоторые клетки остались \textit{липкими}. Один из двух ферзей тоже сразу покрасился, а другой остался липким. \textit{Покрашенный} ферзь может ходить по покрашенным клеткам, а \textit{липкий} ферзь только по липким (оба могут ходить по своим клеткам по несколько раз). Докажите, что какой-то из ферзей может пройтись по всем своим клеткам.

\problem{5}
Найдите все натуральные числа, которые нельзя представить в виде $\dfrac{m}{n} + \dfrac{m+1}{n+1}$, где $m, n$ натуральные числа.

\newpage

\centerline{\bf \Huge Тренировочная олимпиада}
\centerline{\bf \Huge 8 класс}

\problem{1} 
 Докажите, что при $x, y \geq \dfrac{1}{2}$ выполняется неравенство:
 
\[ x^2y^2 + 2(x+y) \geq 4xy + 1 \]



\problem{2}
Остроугольный треугольник $ABC$ с центром описанной окружности в точке $O$ и центром вписанной окружности в точке $I$. Через точку $O$ проведена прямая, перпендикулярная $BI$, а через точку $I$ $-$ прямая, перпендикулярная $BO$. Оказалось, что они пересеклись в точке $A$. Найдите углы треугольника.



\problem{3} 
\textbf{Взлома кода Галактики.} Алина, подружившись с Гроксами попала в их огромную галактическую империю. Все планеты гроксов пронумерованы от 1 до 1000000000. Чтобы добраться от планеты под номером $a$ до планеты под номером $b$ нужно потратить НОД$(a, b)$ спорлингов. Может ли Алина, зная только цены перелётов между планетами, восстановить нумерацию?


\problem{4} 
Три квадратных трехчлена $P(x), Q(x), R(x)$ со старшими коэффициентами, равными 1, имеют по два действительных корня и удовлетворяют условиям 
\[P(Q(x)) = (x-1)(x-3)(x-5)(x-7), \ Q(R(x)) = (x-2)(x-4)(x-6)(x-8) \]

при всех действительных $x$. Чему может быть равно число $P(0)+Q(0)+R(0) $


\problem{5} 
Расул Владимирович составил каждому из 14 учеников индивидуальный вариант тренировочной Олимпиады, но при раздаче каким-то ученикам выдал не свои условия. Назовём обмен \textit{шкибидишным}, если некоторые (не обязательно все) ученики разибиваются на пары и внутри этой пары меняются между собой условиями. Докажите, что за два шкибидишных обмена можно сделать так, что у каждого будет его условие.